\chapter{The probability toolkit}\label{ch:toolkit}

\begin{chapterplan}
Which tools does the paper use? Bayes' rule with likelihood ratios serves the beliefs (Chapter~\ref{ch:beliefs}).
Distribution functions, the tail formula, the win probability and dominance serve the auction
(Chapter~\ref{ch:auction}). Expectations at a belief serve the later chapters.
\end{chapterplan}

The challenger has class \(j\in\{H,L\}\), prior \(\pi=\Pr(j=H)\in(0,1)\), and value \(\theta\) with distribution \(G_j\). The incumbent has value \(R\).

% =====================================================================
\section{Bayes' rule}\label{sec:toolkit-bayes}

Let a signal \(S\) take the value \(s\) with probability \(f_H(s)=\Pr(S=s\mid H)\) in class \(H\) and
\(f_L(s)=\Pr(S=s\mid L)\) in class \(L\). The law of total probability gives \(\Pr(S=s)=\pi f_H(s)+(1-\pi)f_L(s)\).
If \(\Pr(S=s)>0\), the posterior is
\begin{equation}\label{eq:toolkit-bayes}
\mu(s)=\Pr(H\mid S=s)=\frac{\pi f_H(s)}{\pi f_H(s)+(1-\pi)f_L(s)}.
\end{equation}
Divide by \(\pi\) and rearrange: \(\mu(s)>\pi\) exactly when \(f_H(s)>f_L(s)\). A signal counts in favour of the class
that makes it more likely. In the paper the signal is the order \(x\), and condition (E1) of Definition~A.1 is
\eqref{eq:toolkit-bayes} with \(s=x\).

The rule needs \(\Pr(S=s)>0\). In the paper every order has positive probability in both classes.

% =====================================================================
\section{Bayes' rule in odds form}\label{sec:toolkit-odds}

The odds of \(\mu\) are \(\mu/(1-\mu)\), and the likelihood ratio of \(s\) is \(f_H(s)/f_L(s)\). Divide
\eqref{eq:toolkit-bayes} by \(1-\mu(s)=(1-\pi)f_L(s)/[\pi f_H(s)+(1-\pi)f_L(s)]\):
\begin{equation}\label{eq:toolkit-odds}
\frac{\mu(s)}{1-\mu(s)}=\frac{\pi}{1-\pi}\cdot\frac{f_H(s)}{f_L(s)}.
\end{equation}
The map from odds \(o\) to the probability \(o/(1+o)\) is strictly increasing. So a bound on the odds gives a bound
on the probability in the same direction.

\begin{keyidea}
If every likelihood ratio lies between \(1/(1+u)\) and \(1+u\), for some \(u>0\), then every posterior lies between
the bottom and the top posterior:
\begin{equation}\label{eq:toolkit-bounds}
m(\pi)=\frac{\pi}{\pi+(1-\pi)(1+u)}\;\le\;\mu(s)\;\le\;M(\pi)=\frac{\pi(1+u)}{\pi(1+u)+1-\pi}.
\end{equation}
\end{keyidea}

\begin{align*}
\frac{\mu(s)}{1-\mu(s)}&\le\frac{\pi(1+u)}{1-\pi}=:o_{\max}
\why{odds form; likelihood ratio at most \(1+u\)}\\
\mu(s)&\le\frac{o_{\max}}{1+o_{\max}}=\frac{\pi(1+u)}{\pi(1+u)+1-\pi}=M(\pi)
\why{\(o\mapsto o/(1+o)\) increases; cancel \(1-\pi\)}
\end{align*}
The bottom posterior follows in the same way with \(o_{\min}=\pi/[(1-\pi)(1+u)]\). In the paper the bound is a
result, with \(u=\lambda/b\) (Chapter~\ref{ch:beliefs}).

\emph{Averages of posteriors.} An observer who sees only that \(S\in A\), such as a challenger who sees a price that
several orders share, holds the belief \(\Pr(H\mid S\in A)=\sum_{s\in A}w_s\,\mu(s)\), with the weights \(w_s=\Pr(S=s)/\Pr(S\in A)\ge0\) that add up to one.
An average of numbers in \([m(\pi),M(\pi)]\) stays in that interval. So every belief at a price obeys the bounds.

\emph{The mirror identity.} Put \(1-\pi\) in place of \(\pi\) in \(m\):
\begin{equation}\label{eq:toolkit-mirror}
1-M(\pi)=\frac{1-\pi}{\pi(1+u)+1-\pi}=m(1-\pi).
\end{equation}
The function \(m\) is increasing. So the paper's \(\kappa(\pi)=\min\{m(\pi),1-M(\pi)\}\) equals \(m(\min\{\pi,1-\pi\})\).

% =====================================================================
\section{Distribution functions, left limits and atoms}\label{sec:toolkit-cdf}

The distribution function of \(\theta\) is \(G(y)=\Pr(\theta\le y)\). It is nondecreasing and continuous from the
right. Its left limit \(G(p^-)=\lim_{y\uparrow p}G(y)=\Pr(\theta<p)\) excludes the value \(p\), and the jump
\(G(p)-G(p^-)=\Pr(\theta=p)\) is an atom. The complement rule gives
\begin{equation}\label{eq:toolkit-geq}
\Pr(\theta\ge p)=1-G(p^-),\qquad \Pr(\theta>p)=1-G(p).
\end{equation}
A bid equal to the reserve is admissible, so a class-\(j\) value fails the reserve with probability \(G_j(p^-)\). The
class ranking asks for \(G_H(p^-)=G_L(p^-)=0\). The class laws may have atoms; the example uses point masses, with
\(G(y)=\1\{y\ge h\}\). The incumbent law \(F\) is continuous on \([0,\bar r]\) and extended by unity above \(\bar r\).
The uniform law on \([0,r]\) has \(F(y)=\min\{y/r,1\}\) for \(y\ge0\).

\begin{watchout}
\(G(p^-)\) and \(G(p)\) differ by the atom at \(p\). For a point mass at the reserve, \(G(p^-)=0\) but \(G(p)=1\).
\end{watchout}

% =====================================================================
\section{Expectation at a belief}\label{sec:toolkit-expectation}

Expectation is linear, \(\E[aX+bY+c]=a\E[X]+b\E[Y]+c\), with no need for independence, and
\(\E[\1\{A\}]=\Pr(A)\). The class expectation \(\E_j\) uses \(G_j\), and the paper writes
\(t_j=\E_j[t(\theta)]\) and \(g_j=\E_j[g(\theta)]\). Averages go in two stages:
\(\E[\varphi(\theta)]=\pi\E_H[\varphi(\theta)]+(1-\pi)\E_L[\varphi(\theta)]\).

Suppose the signal is independent of the value given the class. Then the signal changes only the weights on the
classes, and with \(\mu=\Pr(H\mid S=s)\)
\begin{equation}\label{eq:toolkit-B}
B(\mu)=\mu\,g_H+(1-\mu)\,g_L=g_L+\mu\,(g_H-g_L).
\end{equation}
Gross profit is a straight line in the belief. If \(g_H>g_L\), then
\(B(\mu)\ge c_H\iff\mu\ge\tau=(c_H-g_L)/(g_H-g_L)\), the preparation threshold. At a fixed entry \(e(x)\), the price is
also a straight line in the posterior: \(P(x)=t_0+e(x)\,[t_L-t_0+\DT\,\mu(x)]\), with \(\DT=t_H-t_L\).

% =====================================================================
\section{The tail formula}\label{sec:toolkit-tail}

\begin{keyidea}
For \(X\ge0\) with distribution function \(G\), \(\E[X]=\int_0^\infty[1-G(y)]\,dy\). Also, for \(c,p\ge0\),
\begin{equation}\label{eq:toolkit-pospart} \E[(X-c)_+]=\int_c^\infty[1-G(y)]\,dy,\qquad \E[\max\{X,p\}]=p+\int_p^\infty[1-G(y)]\,dy. \end{equation}
If \(R\) has a continuous distribution function \(F\) and \(\theta\ge p\) is fixed,
\begin{equation}\label{eq:toolkit-fixed}
\E_R\bigl[(\theta-\max\{p,R\})_+\bigr]=\int_p^\theta F(y)\,dy.
\end{equation}
\end{keyidea}

Each proof has two steps. Write the variable as an integral of an indicator, and then swap the expectation and the
integral. The swap is valid because the integrand is nonnegative (Tonelli). For the first formula,
\(X=\int_0^\infty\1\{y<X\}\,dy\), so \(\E[X]=\int_0^\infty\Pr(X>y)\,dy\).
For \eqref{eq:toolkit-pospart} the integral starts at \(c\). The formula for \(\max\{X,p\}\) uses
\(\max\{X,p\}=p+(X-p)_+\). For \eqref{eq:toolkit-fixed}, check the three cases \(R<p\), \(p\le R<\theta\) and
\(R\ge\theta\) to get \((\theta-\max\{p,R\})_+=\int_p^\theta\1\{R<y\}\,dy\). Then \(\Pr(R<y)=F(y)\).

Formula \eqref{eq:toolkit-fixed} is the challenger's gross profit \(g(\theta)\). The challenger keeps the slice of its value at the level \(y\) exactly when the incumbent's value lies below \(y\), which has probability \(F(y)\).

With the uniform incumbent and a point mass at \(h\), \(p<r<h\):
\(g_H=\int_p^r(y/r)\,dy+\int_r^h1\,dy=h-r/2-p^2/(2r)\), because \(F=1\) above \(r\).

% =====================================================================
\section{Maxima, minima and the probability to win}\label{sec:toolkit-maxmin}

For independent \(X\) and \(Y\), ``both'' becomes a product: \(\Pr(\max\{X,Y\}\le y)=F_X(y)F_Y(y)\) and \(\Pr(\min\{X,Y\}>y)=[1-F_X(y)][1-F_Y(y)]\).
With a continuous \(F\), a fixed value \(\theta\) beats \(R\) with probability \(F(\theta)\), and a tie has probability
zero. A prepared challenger wins when \(\theta\ge p\) and \(R<\theta\), so
\begin{equation}\label{eq:toolkit-win}
W_j(F)=\E_j\bigl[F(\theta)\,\1\{\theta\ge p\}\bigr].
\end{equation}
In the example, \(W_H=F(h)=1\) and \(W_L=\ell/r\), for \(p<\ell<r<h\). Under the class ranking the target receives
\(\max\{p,\min(R,\theta)\}\). The minimum formula and \eqref{eq:toolkit-pospart} then give \(t_j=p+\int_p^\infty[1-F(y)]\,[1-G_j(y)]\,dy\).

% =====================================================================
\section{First-order stochastic dominance}\label{sec:toolkit-fosd}

\(\tilde G\) dominates \(G\) when \(\tilde G(y)\le G(y)\) at every \(y\). The class ranking is \(G_H\le G_L\); a strengthening of the incumbent is \(\tilde F\le F\).
For \(X\ge0\) and \(\varphi'\ge0\), the tail formula extends to
\(\E[\varphi(X)]=\varphi(0)+\int_0^\infty\varphi'(y)[1-G(y)]\,dy\). So \(\E_{\tilde G}[\varphi(X)]-\E_G[\varphi(X)]=\int_0^\infty\varphi'(y)\,[G(y)-\tilde G(y)]\,dy\ge0\).
The class difference \(D(y)=G_L(y)-G_H(y)\ge0\) is the extra probability that a class-\(H\) value passes \(y\), and
\(\E_H[\theta]-\E_L[\theta]=\int_0^\infty D(y)\,dy\). For point masses, \(D(y)=\1\{\ell\le y<h\}\). Uniform incumbents
form a path: for \(r'>r\), \(y/r'\le y/r\) on \([0,r]\) and \(y/r'\le1\) on \((r,r']\).

\begin{watchout}
Dominance means a \emph{smaller} distribution function. A stronger incumbent lowers \(F\), so a challenger at any
fixed value wins less often. Two laws whose distribution functions cross are not ranked.
\end{watchout}

\begin{chaptersummary}
A bound \(1+u\) on every ratio keeps every
posterior, and every average of posteriors, in \([m(\pi),M(\pi)]\). Gross profit at a belief is the straight line
\(B(\mu)=g_L+\mu(g_H-g_L)\). The tail formula gives \(g(\theta)=\int_p^\theta F(y)\,dy\). A stronger incumbent has a smaller \(F\) at every level.
\end{chaptersummary}
